\documentclass[10pt, a5paper]{extarticle}

\usepackage{fontspec}
\usepackage{polyglossia}
\setdefaultlanguage{russian}
\setotherlanguage{english}
\setmainfont{FreeSerif}

\usepackage[
  a5paper,
  top=15mm, bottom=15mm,
  left=20mm, right=20mm
]{geometry}

\usepackage{setspace}
\singlespacing
\setlength{\parindent}{1.25cm}
\setlength{\parskip}{0pt}

\usepackage{amsmath, amssymb}
\pagestyle{empty}
\usepackage[hidelinks]{hyperref}

% Разрешаем чуть большие пробелы ради корректных переносов
\sloppy
% Запрещаем строки с одним словом
\clubpenalty=10000
\widowpenalty=10000

\begin{document}

\noindent УДК <ПОКА ПУСТО>

\vspace{0.25\baselineskip}

\begin{center}
\textbf{ПРОГРАММНЫЙ ИНСТРУМЕНТ\\ ДЛЯ АНАЛИЗА И ГЕНЕРАЦИИ\\
КРИПТОГРАФИЧЕСКИ СТОЙКИХ S-БЛОКОВ}
\end{center}

\vspace{0.25\baselineskip}

\begin{center}
Э.\,Э.\,Морозов
\end{center}

\vspace{0.2\baselineskip}

\begin{center}
Самарский национальный исследовательский университет\\
имени академика С.\,П.\,Королева, г.\,Самара
\end{center}

\vspace{0.25\baselineskip}

\noindent\textbf{Ключевые слова:} S-блок, булевы функции, нелинейность,
дифференциальная равномерность, алгебраическая иммунность,
генерация подстановок, имитация отжига.

\vspace{0.25\baselineskip}

Блочные шифры занимают центральное место в современной симметричной
криптографии, обеспечивая конфиденциальность данных в телекоммуникационных
системах, финансовой инфраструктуре и государственных информационных системах.
Ключевым нелинейным элементом большинства блочных шифров является блок
подстановки (S-блок) --- биективное отображение, определяющее стойкость шифра
к дифференциальному~[1] и линейному~[2] криптоанализу.
Криптографические требования к S-блокам были формализованы
в работе~[3]: нелинейность, дифференциальная равномерность
и алгебраическая степень должны одновременно удовлетворять
установленным пороговым значениям.

Нелинейность $\mathcal{N}(S)$ определяется как минимальное расстояние Хэмминга
между компонентными булевыми функциями S-блока и множеством аффинных функций;
дифференциальная равномерность $\delta(S)$ --- как максимум таблицы разностных
распределений~(DDT); алгебраическая степень $\deg(S)$ --- как максимальная
степень мономов в алгебраической нормальной форме компонентных функций;
алгебраическая иммунность $\mathrm{AI}(S)$ ---как минимальная степень ненулевого аннулятора
компонентных функций.
Между указанными характеристиками существуют теоретические
противоречия: оптимизация одной нередко приводит к ухудшению другой,
что делает проектирование S-блока задачей многокритериальной
оптимизации~[4]. Анализ действующих
стандартов --- AES~[5], «Кузнечик» (ГОСТ Р 34.12-2015)~[6], SM4~[7] и Camellia~[8] --- позволяет установить целевые пороговые значения
для направленного поиска новых подстановок.

Несмотря на то что S-блоки действующих стандартов удовлетворяют
установленным требованиям, ряд теоретических вопросов остаётся открытым.
Точные границы достижимых значений нелинейности для биективных функций
размерности $8\times8$ не установлены: теоретический максимум для
несбалансированных функций равен~120 и достигается бент-функциями,
однако бент-функции не являются биективными и непригодны для
конструирования S-блоков. Наилучшее известное значение нелинейности
для сбалансированного случая составляет~116~[9], тогда как S-блок
AES достигает лишь~112. Кроме того, S-блоки AES и Camellia
обладают компактным алгебраическим описанием через инверсию
в~$\mathrm{GF}(2^8)$, что потенциально облегчает алгебраический анализ.
Отсутствие унифицированного инструмента с полным набором вычисляемых
характеристик затрудняет воспроизводимость криптографических
исследований в данной области.

В рамках настоящей работы разрабатывается инструмент, включающий два
взаимосвязанных модуля. \textit{Модуль анализа} принимает на вход
произвольный S-блок $8\times8$ и вычисляет нелинейность,
дифференциальную равномерность, алгебраическую степень и алгебраическую
иммунность, а также проверяет биективность и отсутствие неподвижных
точек. \textit{Модуль генерации} реализует направленный поиск методом
имитации отжига (\textit{simulated annealing})~[10,11]: операцией мутации
служит случайная транспозиция двух элементов таблицы подстановки,
сохраняющая биективность. Целевая функция --- взвешенная сумма штрафов
за нарушение пороговых значений:
\begin{align*}
  F(S) = \; & w_1\cdot\max(0,\;112-\mathcal{N}(S)) + w_2\cdot\max(0,\;\delta(S)-4)           \\
           + & w_3\cdot\max(0,\;7-\deg(S))             
           + w_4\cdot\max(0,\;\tau^*-\mathrm{AI}(S)),
\end{align*}
где $w_1, w_2, w_3, w_4$ --- настраиваемые веса;\\
\phantom{где} $\tau^*$ --- целевое значение алгебраической иммунности;\\
поиск завершается при $F(S)=0$ либо по достижении заданного числа
итераций.

Планируемые результаты включают верификацию модуля анализа на S-блоках
действующих стандартов; получение подстановок, удовлетворяющих условиям
$\mathcal{N}\geq112$, $\delta\leq4$, $\deg\geq7$, $\mathrm{AI} \geq 3$.

\vspace{0.5\baselineskip}
\noindent\textbf{Список использованных источников}
\vspace{0.2\baselineskip}

\begin{enumerate}
  \setlength{\itemsep}{0pt}
  \setlength{\parsep}{0pt}
  \item Biham~E., Shamir~A. Differential cryptanalysis of DES-like
    cryptosystems~// Journal of Cryptology.~--- 1991.~--- Vol.~4,
    №~1.~--- P.~3--72.
  \item Matsui~M. Linear cryptanalysis method for DES cipher~//
    Advances in Cryptology~--- EUROCRYPT~1993. LNCS~765.~---
    Springer, 1994.~--- P.~386--397.
  \item Carlet~C. Boolean Functions for Cryptography and Coding 
  Theory.~--- Cambridge University Press, 2021.~--- 577~p.
  \item Fuller~J., Millan~W., Dawson~E. Multi-objective optimisation 
  of bijective S-boxes~// New Generation Computing.~--- 
  Springer, 2005.~--- Vol.~23, №~3.~--- P.~201--218.
  \item FIPS~197: Advanced Encryption Standard~(AES).~--- NIST,
    2001.~--- 51~p.
  \item ГОСТ~Р~34.12-2015. Информационная технология.
    Криптографическая защита информации. Блочные шифры.~---
    М.: Стандартинформ, 2015.
  \item GM/T~0002-2012. SM4 Block Cipher Algorithm.~---
    OSCCA, 2012.
  \item Aoki~K., Ichikawa~T., Kanda~M. et~al.
    Camellia: A 128-bit block cipher suitable for multiple
    platforms~// Selected Areas in Cryptography. LNCS~2012.~---
    Springer, 2001.~--- P.~39--56.
  \item Carlet~C., Duong~P.\,P., Dang~T.\,K. et~al.
    Constructing $8\times8$ S-boxes with optimal boolean function
    nonlinearity~// Cryptography.~--- 2025.~--- Vol.~9, №~4.~---
    P.~67. DOI:~10.3390/cryptography9040067.
  \item Kirkpatrick~S., Gelatt~C.\,D., Vecchi~M.\,P. Optimization
    by simulated annealing~// Science.~--- 1983.~--- Vol.~220,
    №~4598.~--- P.~671--680.
    \item Clark~J.,A., Jacob~J.,L., Stepney~S. The design of S-boxes by simulated annealing // New Generation Computing. — 2005. — Vol.~23, №~3. — P.~219–231
\end{enumerate}

\noindent\textit{Морозов Эрнест Эдуардович}, студент кафедры
безопасности информационных систем. E-mail: f3dos1k <at> yandex.ru.

\end{document}